%% ---------------------------------------------------------------------
\section{¿Por qué una teoría de estabilidad?}
\label{sec:ly-porque}
%% ---------------------------------------------------------------------

Un controlador casi nunca se diseña para una sola condición inicial. Se diseña para que el sistema \emph{vuelva} a un punto de operación después de una perturbación (el problema de \clave{regulación}) o para que el error entre la salida y una referencia se haga pequeño (el problema de \clave{seguimiento}). Las dos preguntas se escriben igual: ¿qué les pasa a las soluciones que empiezan \emph{cerca} de un equilibrio? Si el equilibrio de interés es $\eq$, el cambio de variable $z=x-\eq$ lo lleva al origen; en el seguimiento, la ecuación del error hace lo mismo. Por eso, en todo el capítulo,
\begin{equation}
  \dot x=f(x),\qquad f:\D\to\R^n\ \text{localmente Lipschitz},\qquad 0\in\D,\qquad f(0)=0 ,
  \label{eq:ly-sistema}
\end{equation}
con $\D\subset\R^n$ un dominio (abierto y conexo, \cref{def:pm-convexo}) y $\flujo{t}{x_0}$ la solución que pasa por $x_0$ en $t=0$. El sistema es \emph{invariante en el tiempo}: $f$ no depende de $t$. Los sistemas variantes en el tiempo, donde las definiciones necesitan uniformidad, son el tema del capítulo~9.

\notamargen{\textbf{Ruta de las clases.} Semana~8: \S\ref{sec:ly-definiciones}--\ref{sec:ly-directo}. Semana~9: \S\ref{sec:ly-global}--\ref{sec:ly-region}. La prueba del \cref{thm:ly-directo} se hace en clase; las demás se leen por cuenta propia.}

En los capítulos~\ref{cap:unidimensionales} y~\ref{cap:orbitas} la estabilidad se juzgó \emph{mirando} el retrato de fase o linealizando. Ninguno de los dos sirve en general: en dimensión mayor que dos no hay retrato que mirar, y la linealización falla justo en los casos interesantes (\cref{sec:ly-indirecto}). Lyapunov propuso en 1892 una salida que es la idea central del curso: buscar una función $V(x)$ que haga de \emph{energía} y probar que decrece a lo largo de las soluciones \emph{sin calcular las soluciones}.

\begin{historia}[De Járkov a la ingeniería de control]
Aleksandr Mijáilovich Lyapunov (1857--1918), alumno de Chebyshov en San Petersburgo, defendió en 1892 su tesis doctoral \emph{El problema general de la estabilidad del movimiento}, escrita en Járkov. Contiene los dos métodos de este capítulo: el \emph{primero} (indirecto), que estudia la ecuación linealizada, y el \emph{segundo} (directo), el de las funciones $V$. La tesis se tradujo al francés en 1907 y fue conocida sobre todo por los matemáticos rusos durante cuarenta años; la traducción completa al inglés apareció en 1992~\citep{Lya92}. Barbashin y Krasovskii dieron en 1952 el resultado global de la \cref{sec:ly-global}~\citep{BaKr52}; LaSalle, en 1960, el principio de invarianza del capítulo~8~\citep{LaS60}. El mismo año, Kalman y Bertram mostraron a los ingenieros de control que el segundo método era una herramienta de \emph{diseño}~\citep{KaBe60}, y desde entonces casi todo el control no lineal se escribe en su lenguaje.
\end{historia}

\begin{figure}[ht]
\centering
\includegraphics[width=0.92\linewidth]{cap07/t_mapa_lyapunov}
\caption{Mapa del capítulo. Los dos métodos de Lyapunov y su punto de encuentro, los sistemas lineales. Lo que ninguno de los dos decide se retoma en el capítulo~8.}
\label{fig:ly-mapa}
\end{figure}

\begin{pausa}
Antes de empezar, apueste: si todas las soluciones de un sistema convergen a un equilibrio, ¿es ese equilibrio estable?
\respuesta{No necesariamente. ``Estable'' pide que las soluciones que empiezan cerca \emph{se queden} cerca; converger después de una excursión grande no basta. El sistema de la \cref{tr:ly-convergencia} tiene todas sus soluciones convergiendo a $(1,0)$, y el equilibrio es inestable.}
\end{pausa}
